\section{Resultados}\label{sec:resultados}

\subsection{Posición inicial}
Evaluada con el modelo de riesgo de la Sección~\ref{sec:riesgo}, la posición inicial rinde en promedio
$\E[\Delta PN] = \PEN3.96$ mm a un año, con una volatilidad de \PEN25.4 mm y un $\CVaR_{95\%}$ de \PEN50.1 mm
(Cuadro~\ref{tab:riesgo-inicial}): una cuarta parte del patrimonio inicial en el 5\,\% de escenarios más
adversos. Un 25\,\% de esa cola de pérdida proviene de los cuatro escenarios de estrés pese a que representan
solo 1.25\,\% de la probabilidad total, y el escenario de mayor pérdida es una depreciación del sol acompañada
de una caída de tasas y de renta variable, que cuesta \PEN45.6 mm. La única restricción del Cuadro~4 del caso
que la posición inicial incumple es el peso máximo de renta variable en los activos (22\,\% contra 20\,\%).

\begin{table}[htbp]
\centering
\caption{Riesgo a un año de la posición inicial (S/ mm salvo probabilidades).}
\label{tab:riesgo-inicial}
\begin{tabular}{lr}
\toprule
 & valor \\
\midrule
PN0 & 200.00 \\
E[ΔPN] & 3.96 \\
E[PN\_T] & 203.96 \\
σ(ΔPN) & 25.43 \\
VaR 95\% & 42.49 \\
CVaR 95\% & 50.12 \\
P(PN\_T < 120) & 0.00 \\
prob. cola: estrés / total & 0.25 \\
ΔPN C\_STRESS\_1 & -1.03 \\
ΔPN C\_STRESS\_2 & -45.61 \\
ΔPN C\_STRESS\_3 & -23.34 \\
ΔPN C\_STRESS\_4 & -8.50 \\
\bottomrule
\end{tabular}
\end{table}


\subsection{Cartera óptima}
Con la aversión al riesgo de referencia del caso, $\lambda = 0.25$, la reasignación óptima
(Cuadro~\ref{tab:posiciones-optimas} y Figura~\ref{fig:posiciones}) concentra los activos en soles: la caja
sube a \PEN200 mm (el máximo de 20\,\% de los activos, no el piso de \PEN80 mm, que queda muy por debajo) y el
bono fijo en soles A02 sube a su tope individual de \PEN350 mm. Del lado en dólares, el bono flotante A05 casi
se duplica a \PEN240 mm mientras que el bono fijo A04 se liquida por completo; la renta variable A06 cae a su
piso implícito de 5\,\% (\PEN50 mm), el mínimo compatible con los máximos de caja y renta fija. En los pasivos,
el préstamo corto en soles L01 sube a su tope de concentración (\PEN280 mm) y el bono fijo L02 baja a
\PEN40 mm, mientras que la mezcla de tasa fija y flotante de los pasivos queda exactamente en sus dos límites
(45\,\% fijo, 55\,\% flotante): son restricciones gemelas, extremos opuestos de la misma partición.

\begin{table}[htbp]
\centering
\caption{Posiciones inicial, óptima (λ = 0.25) y alternativas (S/ mm).}
\label{tab:posiciones-optimas}
\resizebox{\ifdim\width>\textwidth\textwidth\else\width\fi}{!}{%
\begin{tabular}{llllrrrrr}
\toprule
 & Lado & Instrumento & Moneda & Inicial & Óptimo (λ = 0.25) & Cambio & Media-varianza & Vencimientos en t\_H (ENFORCE) \\
\midrule
A01 & Activo & Caja PEN & PEN & 150.0 & 200.0 & 50.0 & 194.7 & 200.0 \\
A02 & Activo & Bono PEN Fijo 2030 & PEN & 180.0 & 350.0 & 170.0 & 350.0 & 350.0 \\
A03 & Activo & Bono PEN Flotante 2031 & PEN & 160.0 & 160.0 & 0.0 & 200.1 & 160.0 \\
A04 & Activo & Bono USD Fijo 2032 & USD & 150.0 & 0.0 & -150.0 & 0.0 & 0.0 \\
A05 & Activo & FRN USD 2030 & USD & 140.0 & 240.0 & 100.0 & 199.9 & 240.0 \\
A06 & Activo & ETF renta variable global & USD & 220.0 & 50.0 & -170.0 & 55.3 & 50.0 \\
L01 & Pasivo & Préstamo PEN CP flotante & PEN & 120.0 & 280.0 & 160.0 & 273.1 & 200.0 \\
L02 & Pasivo & Bono PEN Fijo 2029 & PEN & 150.0 & 40.0 & -110.0 & 40.0 & 103.7 \\
L03 & Pasivo & Bono PEN Fijo 2036 & PEN & 180.0 & 180.0 & 0.0 & 180.0 & 158.2 \\
L04 & Pasivo & Préstamo PEN Flotante 2032 & PEN & 110.0 & 110.0 & 0.0 & 110.0 & 110.0 \\
L05 & Pasivo & Bono USD Fijo 2031 & USD & 140.0 & 140.0 & 0.0 & 140.0 & 140.0 \\
L06 & Pasivo & Préstamo USD Flotante 2035 & USD & 100.0 & 50.0 & -50.0 & 56.9 & 88.0 \\
\bottomrule
\end{tabular}}
\end{table}


\begin{figure}[htbp]\centering
  \caption{Posiciones inicial y óptima por instrumento (barras) y de la variante con vencimientos exigidos en
  $t_H$ (marca horizontal).}\label{fig:posiciones}
  \includegraphics[width=\linewidth]{posiciones_inicial_optima.pdf}
\end{figure}

Esa reasignación paga \PEN2.73 mm en costos de transacción y sube el resultado esperado bruto a
\PEN5.83 mm; aun así, el resultado neto (\PEN3.11 mm) queda por debajo del de la posición inicial
(\PEN3.96 mm). El modelo cede \PEN0.85 mm de valor esperado a cambio de bajar la volatilidad del patrimonio
de \PEN25.4 mm a \PEN7.0 mm y el $\CVaR_{95\%}$ de \PEN50.1 mm a \PEN7.8 mm, una caída de 84\,\%
(Cuadro~\ref{tab:metricas-optimizacion}). La rentabilidad esperada de los activos apenas cambia (de 6.08\,\% a
6.01\,\%), pero el costo esperado de los pasivos baja de 7.10\,\% a 6.79\,\%: la ganancia de la optimización
viene sobre todo de una estructura de pasivos más barata y menos expuesta a la cola, no de activos más
rentables.

\begin{table}[htbp]
\centering
\caption{Métricas a un año (S/ mm salvo indicación). α = 0.95; objetivo evaluado con λ = 0.25.}
\label{tab:metricas-optimizacion}
\resizebox{\ifdim\width>\textwidth\textwidth\else\width\fi}{!}{%
\begin{tabular}{lrrrr}
\toprule
 & Inicial & Óptimo (λ = 0.25) & Media-varianza & Venc. exigidos al horizonte \\
\midrule
Objetivo (E − λ·CVaR − TC) & -9.34 & 0.91 & 0.70 & 0.70 \\
E[ΔPN] & 3.29 & 5.56 & 5.55 & 4.53 \\
Costo de transacción & 0.00 & 2.72 & 2.71 & 1.83 \\
E[ΔPN] − TC & 3.29 & 2.84 & 2.84 & 2.70 \\
E[PN] al horizonte & 203.29 & 202.84 & 202.84 & 202.70 \\
σ(ΔPN) & 25.33 & 6.27 & 5.74 & 5.66 \\
VaR 95 \% & 43.07 & 3.88 & 3.06 & 2.78 \\
CVaR 95 \% & 50.52 & 7.71 & 8.55 & 8.02 \\
Rentabilidad esperada activos (\%) & 5.95 & 5.94 & 5.94 & 5.94 \\
Costo esperado pasivos (\%) & 7.02 & 6.73 & 6.73 & 6.86 \\
Estrés en la cola (\%) & 25.00 & 25.00 & 25.00 & 25.00 \\
Pérdida C\_STRESS\_1 & 1.03 & -13.30 & -5.07 & -11.84 \\
Pérdida C\_STRESS\_2 & 45.61 & 13.90 & 19.85 & 20.35 \\
Pérdida C\_STRESS\_3 & 23.34 & -13.02 & -2.21 & -6.51 \\
Pérdida C\_STRESS\_4 & 8.50 & -19.84 & -12.75 & -21.99 \\
\bottomrule
\end{tabular}}
\end{table}


\subsection{Barrido de la aversión al riesgo}
Variar $\lambda$ entre 0 y 5 traza la frontera entre resultado esperado neto y $\CVaR_{95\%}$
(Cuadro~\ref{tab:barrido-lambda} y Figura~\ref{fig:frontera}). Con $\lambda = 0$ el objetivo ignora el riesgo y
elige una solución de esquina: 49\,\% de activos en dólares y 20\,\% en renta variable, con un
$\CVaR_{95\%}$ de \PEN61.0 mm, más riesgosa que la propia posición inicial. Subir $\lambda$ a 0.1 ya recorta
ese CVaR a \PEN32.1 mm, y en $\lambda = 0.25$ la solución colapsa a los límites de caja y renta fija descritos
arriba. Para $\lambda \ge 0.5$ el resultado esperado neto cae por debajo de \PEN1.1 mm y en $\lambda = 5$ se
vuelve negativo (\mbox{$-$\PEN0.69} mm): a partir de ahí, evitar más cola cuesta más de lo que queda de
resultado esperado. La composición de la cartera se mueve en la misma dirección: la renta variable salta del
20\,\% al piso de 5\,\% no bien $\lambda$ deja de ser cero, y el peso de pasivos en dólares sube de 23.75\,\%
a 43\,\% a medida que $\lambda$ crece, porque financiarse en dólares resulta, en este conjunto de escenarios,
más barato en la cola que financiarse en soles.

\begin{table}[htbp]
\centering
\caption{Barrido de λ: resultado esperado, costo, riesgo y composición (S/ mm y \% del lado).}
\label{tab:barrido-lambda}
\resizebox{\ifdim\width>\textwidth\textwidth\else\width\fi}{!}{%
\begin{tabular}{lrrrrrrrrr}
\toprule
 & E[ΔPN] & TC & E[ΔPN] − TC & σ(ΔPN) & CVaR 95 \% & Activos USD (\%) & Equity (\%) & Pasivos USD (\%) & Pasivos flotantes (\%) \\
\midrule
0.000000 & 9.15 & 2.31 & 6.84 & 30.22 & 60.97 & 49.00 & 20.00 & 23.75 & 55.00 \\
0.100000 & 8.26 & 2.17 & 6.09 & 19.79 & 32.06 & 41.00 & 17.00 & 20.00 & 47.60 \\
0.250000 & 5.83 & 2.73 & 3.11 & 7.03 & 7.82 & 29.00 & 5.00 & 23.75 & 55.00 \\
0.500000 & 7.26 & 6.13 & 1.12 & 5.96 & 2.54 & 27.46 & 5.00 & 27.63 & 52.13 \\
1.000000 & 7.41 & 6.62 & 0.79 & 6.09 & 2.01 & 28.35 & 5.00 & 28.56 & 51.06 \\
2.000000 & 7.62 & 7.61 & 0.01 & 5.91 & 1.52 & 33.73 & 5.00 & 35.37 & 53.45 \\
5.000000 & 7.96 & 8.66 & -0.69 & 6.02 & 1.35 & 40.00 & 5.00 & 43.11 & 55.00 \\
\bottomrule
\end{tabular}}
\end{table}


\begin{figure}[htbp]\centering
  \caption{Frontera resultado esperado neto vs. $\CVaR_{95\%}$ sobre el barrido de $\lambda$, con la posición
  inicial, el punto de comparación de mínima varianza y la variante con vencimientos exigidos en $t_H$.}
  \label{fig:frontera}
  \includegraphics[width=0.65\linewidth]{frontera_lambda.pdf}
\end{figure}

\subsection{Comparación con mínima varianza y con vencimientos exigidos en el horizonte}
Fijando el mismo resultado esperado neto que la cartera óptima (\PEN3.11 mm), una cartera de mínima varianza
con las mismas restricciones logra una volatilidad algo menor (\PEN6.40 mm contra \PEN7.03 mm) pero un
$\CVaR_{95\%}$ mayor (\PEN9.23 mm contra \PEN7.82 mm): optimizar directamente la medida de riesgo que exige el
caso controla mejor la cola que minimizar la dispersión total, aunque acepte algo más de varianza en el
centro de la distribución. Exigir que los límites de vencimiento del Cuadro~4 del caso se cumplan también en $t_H$, y no solo en $t_0$ como
hace la solución base, cuesta \PEN0.29 mm de objetivo (de 1.16 a 0.87). El préstamo corto L01 baja de
\PEN280 mm a \PEN200 mm y el $\CVaR_{95\%}$ mejora a \PEN7.33 mm, porque esa restricción adicional también
limita, de paso, la concentración de la cartera óptima base.

\begin{figure}[htbp]\centering
  \caption{Distribución de $\Delta PN$ a un año: posición inicial, cartera óptima y mínima varianza
  (histograma: escenarios de remuestreo; triángulos: escenarios de estrés; línea punteada: $-\CVaR_{95\%}$ de
  cada cartera).}\label{fig:distribucion-optima}
  \includegraphics[width=0.75\linewidth]{distribucion_dpn_optima.pdf}
\end{figure}

La Figura~\ref{fig:distribucion-optima} resume el efecto conjunto: la distribución de la cartera óptima es
mucho más angosta que la inicial y sus escenarios de estrés, que en la posición inicial caían muy por debajo
del grueso de la distribución, quedan cerca del cuerpo de la distribución remuestreada.
